\section{Learning curves and detailed test results}
\label{app:figures}
\begin{figure}[H]
\centering
\includegraphics[width=\textwidth]{baseline_curves.pdf}
\caption{Baseline learning curves. Dashed grey: sampled training loss; solid colors: sampled validation loss. Stars mark best recorded validation checkpoints. The bottom row omits initialization. Lines connect recorded checks; they do not add observations.}
\label{fig:baseline}
\end{figure}
\begin{figure}[H]
\centering
\includegraphics[width=\textwidth]{dropout_comparison.pdf}
\caption{Follow-up prompted by the baseline curves: dropout 0.1 (grey dashed) versus 0.2 (colored solid). Each size is compared with its own baseline; stars mark its best recorded validation losses. Medium and large improve, while small does not.}
\label{fig:dropout}
\end{figure}
\clearpage
\begin{figure}[H]
\centering
\includegraphics[width=0.70\textwidth]{test_piece_distribution.pdf}
\caption{Test-loss distributions for the selected settings: small dropout 0.1, medium/large dropout 0.2. Each model is scored on the same 177 pieces. These distributions give each piece equal weight, while the overall loss gives each token equal weight. They do not measure variation across training seeds.}
\label{fig:pieceevaluation}
\end{figure}
\begin{table}[H]
\centering
\caption{Exploratory checks on the same saved test predictions. Equal-piece loss gives every piece the same weight. The last column counts pieces where small has lower loss. These are not repeated training runs.}
\label{tab:pairedtest}
\begin{tabular}{lrrr}
\toprule Model & Token-weighted loss & Equal-piece loss & Small wins / 177\\\midrule
Small & 1.6935 & 1.7090 & --\\
Medium & 1.7179 & 1.7335 & 177\\
Large & 1.7277 & 1.7432 & 177\\
\bottomrule
\end{tabular}
\end{table}

\begin{table}[H]
\centering
\caption{Per-event-family test cross-entropy for the selected configurations. Counts are shared across models.}
\label{tab:eventdiagnostics}
\begin{tabular}{lrrrr}
\toprule Family & Targets & Small ($p=0.1$) & Medium ($p=0.2$) & Large ($p=0.2$)\\\midrule
NOTE\_ON & 741,410 & 1.7608 & 1.8615 & 1.8960\\
NOTE\_OFF & 741,410 & 1.1307 & 1.1544 & 1.1721\\
TIME\_SHIFT & 1,999,264 & 1.6817 & 1.6804 & 1.6805\\
VELOCITY & 741,410 & 2.2357 & 2.2538 & 2.2692\\
PEDAL & 110,421 & 1.5642 & 1.5862 & 1.5194\\
EOS & 177 & 20.5116 & 21.6459 & 20.2792\\
\bottomrule
\end{tabular}
\end{table}

The event counts are identical across models. Each category's loss still uses probabilities over the full vocabulary. EOS has high loss but only 177 targets, so poor prediction of endings has little effect on overall loss. These diagnostics do not measure audio or listening quality. Additional figures and analysis notebooks remain available in the repository.
